% =====================================================================
%  Part III -- The reaction network as a formal object.
%  Source: tier0.md Part II (1442-1828) then 0.5 (643-757),
%          0.6 (759-844), 0.7 (846-946).
%  Ordered so that every pointer from the old 0.5-0.7 into Part II
%  resolves backward.
%  Self-check merges II.8 (minus items 4 and 7, relocated to Part IV)
%  with 0.7.1.
% =====================================================================
\part{The reaction network as a formal object}
\label{part:network}

% ---------------------------------------------------------------------
\chapter{Nuclide bookkeeping}
\label{sec:bookkeeping}

A nuclide is specified by $(Z, N)$: $Z$ protons, $N$ neutrons, $A = Z + N$.

\begin{itemize}
\item Names are \textbf{MESA/pynucastro chem ids}, lowercase: \code{si28},
      \code{fe56}, \code{he4}, and critically \code{neut} for the free neutron,
      \code{h1} for the free proton.
\item Free nucleons are \textbf{network species}, not a background bath.
      \code{neut} carries ($Z=0$, $A=1$); \code{h1} carries ($Z=1$, $A=1$). At
      silicon-burning temperatures photodisintegration liberates free nucleons
      in abundance --- the Sobol initial compositions carry median
      $X_{\mathrm{neut}} \approx 1.7\times10^{-2}$. They participate in the
      same algebra as every other species, and this is not incidental:
      \S\ref{sec:lastday} established that Si burning \emph{is} the traffic in
      free nucleons and $\alpha$ particles.
\end{itemize}

% ---------------------------------------------------------------------
\chapter{Binding energy and why the iron peak exists}
\label{sec:bindingenergy}

\begin{equation*}
B(Z,N) = \bigl[Z m_p + N m_n - m(Z,N)\bigr]\,c^2
\end{equation*}

Semi-empirical mass formula \citep[][this coefficient set is the least-squares
fit of \citealt{eisberg1985}]{krane1988}:

\begin{equation*}
B \approx a_V A - a_S A^{2/3} - a_C \frac{Z(Z-1)}{A^{1/3}}
        - a_A \frac{(A-2Z)^2}{A} \pm \delta
\end{equation*}

\begin{center}
\small
\begin{tabularx}{\textwidth}{@{}l l Y@{}}
\toprule
\textbf{Term} & \textbf{Coefficient} & \textbf{Effect} \\
\midrule
Volume    & $a_V \approx 15.8$~MeV & grows as $A$ --- favours heavy nuclei \\
Surface   & $a_S \approx 18.3$     & penalizes small $A$ --- $B/A$ rises
                                     steeply at low $A$ \\
Coulomb   & $a_C \approx 0.714$    & grows as $Z^2/A^{1/3}$ --- penalizes heavy
                                     nuclei \\
Asymmetry & $a_A \approx 23.2$     & penalizes $N \neq Z$ \\
Pairing   & $\pm\delta$            & even--even favoured \\
\bottomrule
\end{tabularx}
\end{center}

Surface losses fade with $A$ while Coulomb losses grow, so $B/A$ peaks near
$A \approx 56$--62 and declines. Fusion is exothermic below the peak,
endothermic above. \textbf{This is why silicon burning is the last stage.}

\begin{warn}
\textbf{The mass-excess trap.} Data tables give the mass \emph{excess} $\Delta$,
and masses are \emph{atomic}. Converting takes
$B = Z\Delta(\nuc{H}{1}) + N\Delta(n) - \Delta$, and $Q$-values are
$Q = \sum\Delta_{\mathrm{in}} - \sum\Delta_{\mathrm{out}}$. Worked examples and
a verification exercise against \code{configs/reactions\_*\_mesa.yaml} are in
\S\ref{sec:massexcess}.
\end{warn}

\section{The subtlety to internalize now}
\label{sec:constrainedoptimum}

The most bound nuclide is \nuc{Ni}{62} ($B/A = 8.795$~MeV;
\citealt{fewell1995}), the most bound at $A = 56$ is \nuc{Fe}{56} (8.790) ---
yet silicon burning at $\Ye \approx 0.5$ produces \textbf{\nuc{Ni}{56}}
(8.643). The reason: the equilibrium composition maximizes binding
\emph{subject to the \Ye{} constraint}. At $\Ye = 0.5$ matter cannot reach
\nuc{Fe}{56} ($\Ye = 26/56 = 0.464$) without weak reactions, which are slow.
\nuc{Ni}{56} is the most bound nuclide with $Z = N$, and doubly magic on top of
that (\S\ref{sec:magicnumbers}).

\textbf{Carry this forward:} ``the equilibrium composition'' is always \emph{at
a given \Ye{}}. That is why the NSE solver has exactly two constraints,
$\sum X = 1$ and \Ye{}, and why the label-pathology finding is stated as a
distance from NSE \emph{at the label's own \Ye{}}.

% ---------------------------------------------------------------------
\chapter{From number density to molar abundance --- derive
  \texorpdfstring{$Y = X/A$}{Y = X/A}}
\label{sec:molarabundance}

The mass density in species $i$ is $\rho X_i$, and each nucleus has mass
$\approx A_i m_u$:

\begin{equation*}
n_i = \frac{\rho X_i}{A_i m_u} = \frac{\rho N_A X_i}{A_i}
\end{equation*}

using $m_u = 1/N_A$ in cgs. Define the \textbf{molar abundance}:

\begin{equation}
Y_i \equiv \frac{X_i}{A_i} \qquad \text{so} \qquad n_i = \rho N_A Y_i
\label{eq:molar-abundance}
\end{equation}

The mass-fraction normalization becomes

\begin{equation}
\sum_i X_i = 1 \qquad \iff \qquad \sum_i A_i Y_i = 1
\label{eq:baryon-conservation}
\end{equation}

Equation \eqref{eq:baryon-conservation} is the baryon conservation law the
entire conservation layer is built on. It is \textbf{linear in $Y$} --- not an
accident; it is the reason $Y$ is the right variable (and the standard network
variable: \citealt{hix1999b} introduce $Y_i = n_i/\rho N_A$ for exactly these
reasons). A pedantic aside on $m_u = 1/N_A$: exact in the pre-2019 cgs
convention; since the 2019 SI redefinition,
$m_u N_A = 1~\mathrm{g/mol} \times (1 + 1.05\times10^{-9})$
\citep{codata2022} --- six orders below anything tabulated here.

\textbf{A precision note.} Strictly $A_i m_u$ is not the nuclear mass. The
deviation is exactly the mass excess of \S\ref{sec:massexcess},
$(M - A u)/(A u) = \Delta/(A \cdot 931.494~\mathrm{MeV})$, and it is worth
tabulating because its \emph{sign and size are both counter-intuitive}:

\begin{center}
\begin{tabular}{@{}l c c c c c@{}}
\toprule
 & \code{neut} & \code{h1} & \code{he4} & \code{si28} & \code{fe56} \\
\midrule
$(M - A u)/(A u)$ & \textbf{$+0.87\%$} & \textbf{$+0.78\%$} & $+0.065\%$
                  & $-0.082\%$ & $-0.116\%$ \\
\bottomrule
\end{tabular}
\end{center}

Two traps here (values from AME2020; note these are atomic-convention numbers
--- the \code{h1} entry includes the electron, and the bare proton would read
$+0.73\%$). First, the largest deviations belong to the \textbf{free nucleons},
and they are \textbf{positive} --- $n$ and $p$ are \emph{heavier} than
$A\cdot u$, because the atomic mass unit is defined by \nuc{C}{12}, which is
itself bound. Second, for the heavy nuclei that carry the mass the deviation is
only ${\sim}10^{-3}$, an order of magnitude smaller. The ${\sim}0.9\%$ figure
that is easy to reach for is $B/(A m_u c^2)$ --- the defect relative to
\emph{free constituent nucleons} --- which is a different quantity and answers
a different question.

Using the integer $A$ means $\sum A_i Y_i = 1$ holds \emph{exactly by
definition} rather than approximately, and the mass defect is accounted
separately in the energy equation ($\enuc = -N_A \sum m_i c^2 \Delta Y_i$, with
the mass excess usable in place of $m_i c^2$ because nucleon number is
conserved --- \citealt{hix1999b}, eq.~12). This is a deliberate bookkeeping
split: \textbf{exact integer arithmetic for the constraint, mass excesses for
the energy.} It is what makes $A\cdot\nu = 0$ hold to machine zero rather than
to the ${\sim}10^{-3}$ set by the heavy-nucleus column of the table above.

\begin{repopointer}
\textbf{In the code:} \code{src/gnn\_nucleo/data/subsample.py::compute\_ye\_initial}
is \code{X @ (Z/A)}, i.e.\ exactly $\sum_i Z_i X_i/A_i$ --- a one-line
implementation of \eqref{eq:molar-abundance} and \eqref{eq:ye}.
\end{repopointer}

% ---------------------------------------------------------------------
\chapter{Why networks integrate \texorpdfstring{$Y$}{Y} and not
  \texorpdfstring{$X$}{X} --- the real reason}
\label{sec:whyY}

\section{Argument 1 --- \texorpdfstring{$Y$}{Y} is a per-baryon quantity,
  invariant under compression}
\label{sec:perbaryon}

From \eqref{eq:molar-abundance}, $Y_i = n_i/(\rho N_A) =$ (number of species
$i$) / (number of baryons). Compressing or expanding a fluid element changes
$\rho$ and every $n_i$ in proportion, leaving every $Y_i$ untouched. So in the
Lagrangian frame:

\begin{equation*}
\left.\frac{\dd Y_i}{\dd t}\right|_{\text{total}}
= \left.\frac{\dd Y_i}{\dd t}\right|_{\text{reactions}}
\end{equation*}

with no advection or compression term. Working in $n_i$ you would carry a
spurious $-n_i(\dot\rho/\rho)$ term in every equation; in $X$ you would carry it
too, hidden inside a nonlinear rate expression.

\section{Argument 2 --- the constraint stays linear with constant coefficients}
\label{sec:linearconstraint}

\begin{equation*}
\begin{aligned}
\text{In } Y: &\quad \sum_i A_i Y_i = 1
  && \text{coefficients } A_i \text{ are integers, fixed forever}\\
\text{In } n: &\quad \sum_i A_i n_i = \rho N_A
  && \text{right-hand side is time-dependent}
\end{aligned}
\end{equation*}

This is what makes the conservation layer a \emph{fixed} linear operator,
precomputed once and applied in decode. In any other variable the constraint
matrix would be state-dependent and the architecture would collapse. Trace the
dependency: linear constraint $\to$ fixed $\nu$ and $C$ $\to$ exact
$A\cdot\nu = 0$ $\to$ conservation survives a random $\varphi$ $\to$
conservation-by-construction is possible at all.

\section{The network ODE}
\label{sec:networkode}

With rate $R_j$ [mol\,g$^{-1}$\,s$^{-1}$] and $\nu_{ij}$ the net stoichiometric
coefficient (products $+$, reactants $-$):

\begin{equation}
\frac{\dd Y_i}{\dd t} = \sum_j \nu_{ij} R_j \qquad \iff \qquad
\frac{\dd\mathbf{Y}}{\dd t} = \nu\,\mathbf{R}
\label{eq:network-ode}
\end{equation}

For a two-body reaction $1 + 2 \to 3$,
$r = n_1n_2\langle\sigma v\rangle/(1+\delta_{12})$ per unit volume; converting
with \eqref{eq:molar-abundance} and $R = r/(\rho N_A)$:

\begin{equation*}
R = \frac{\rho N_A Y_1 Y_2 \langle\sigma v\rangle}{1+\delta_{12}}
\end{equation*}

--- one net power of $N_A$ (\citealt{hix1999b}, eq.~10: two-body terms carry
$\rho N_A$, three-body $\rho^2 N_A^2$); equivalently
$R = \rho Y_1 Y_2 \cdot [N_A\langle\sigma v\rangle]/(1+\delta_{12})$, where
$N_A\langle\sigma v\rangle$ is the tabulated REACLIB quantity.

The general form --- density exponent $\rho^{N-1}$ for $N$ reactants, $1/n!$
for $n$ identical reactants --- is derived in S2. Note now only that $R$ is a
\textbf{product of abundances times a function of $(T, \rho)$}, which makes
$\partial R/\partial Y$ analytically available by the product rule --- the fact
S10's integrator Jacobian rests on.

% ---------------------------------------------------------------------
\chapter{The electron fraction \texorpdfstring{\Ye{}}{Ye} --- the central
  variable}
\label{sec:ye}

\section{Derivation}
\label{sec:yederivation}

Charge neutrality: electron number density equals total proton number density,

\begin{equation*}
n_e = \sum_i Z_i n_i = \rho N_A \sum_i Z_i Y_i
\end{equation*}

Define, in exact parallel to \eqref{eq:molar-abundance}:

\begin{equation}
Y_e \equiv \frac{n_e}{\rho N_A} = \sum_i Z_i Y_i = \sum_i \frac{Z_i X_i}{A_i}
\label{eq:ye}
\end{equation}

\Ye{} is the number of electrons per baryon. For any nuclide with $Z = A/2$,
$Z_i/A_i = 1/2$, so symmetric matter has $\Ye = 0.5$ exactly. $\Ye < 0.5$ is
neutron-rich.

\section{The neutron excess}
\label{sec:neutronexcess}

\begin{equation}
\eta = \sum_i (N_i - Z_i)\,Y_i
     = \sum_i A_i Y_i - 2\sum_i Z_i Y_i = 1 - 2Y_e
\label{eq:neutron-excess}
\end{equation}

So $\eta$ and \Ye{} carry the same information ($\eta$ is the standard neutron
excess --- \citealt{hix1999b}). The literature uses both, and confusing them is
a real source of misread thresholds --- this repo carries an explicit open item
on exactly that ambiguity in a cited paper.

\textbf{Percentages quoted ``in $\eta$'' and ``in \Ye{}'' are not
interchangeable}, and the amplification is larger than it first looks.
Differentiating \eqref{eq:neutron-excess} gives $\delta\eta = \mathbf{-2}\,
\delta\Ye$, so the relative-error amplification is

\begin{equation*}
\frac{|\delta\eta|/\eta}{|\delta Y_e|/Y_e}
  \;=\; \frac{2Y_e}{\eta} \;=\; \frac{2Y_e}{1-2Y_e}
\end{equation*}

At $\Ye = 0.47$, $\eta = 0.06$ and the factor is
$2(0.47)/0.06 = \mathbf{15.7}$: a 1\% error in \Ye{} is a
\textbf{${\sim}16\%$} error in $\eta$, not 8\%. (Dropping the factor of 2 from
\eqref{eq:neutron-excess} is the easy way to get 7.8 and it is wrong.) The
amplification diverges as $\Ye \to 0.5$ --- precisely the upper edge of the
regime box --- which is why $\eta$ is the worse variable to write a tolerance
in here, and why every gate in this project is stated in \Ye{}.

\section{What changes \texorpdfstring{\Ye{}}{Ye}}
\label{sec:whatchangesye}

Strong and electromagnetic reactions conserve $Z$ and $A$ separately, so they
cannot change \Ye{} at all. Only \textbf{weak reactions} can:

\begin{center}
\small
\begin{tabularx}{\textwidth}{@{}R{1.000} l l l@{}}
\toprule
\textbf{Process} & \textbf{Effect on $(Z, N)$} & \textbf{Effect on \Ye{}} &
\textbf{Emits} \\
\midrule
electron capture, $e^- + p \to n + \nu$   & $Z-1$, $N+1$ &
  \textbf{decreases} & $\nu$ \\
$\beta^+$ decay, $p \to n + e^+ + \nu$    & $Z-1$, $N+1$ &
  \textbf{decreases} & $\nu$ \\
$\beta^-$ decay, $n \to p + e^- + \bar\nu$ & $Z+1$, $N-1$ &
  \textbf{increases} & $\bar\nu$ \\
positron capture, $e^+ + n \to p + \bar\nu$ & $Z+1$, $N-1$ &
  \textbf{increases} & $\bar\nu$ \\
\bottomrule
\end{tabularx}
\end{center}

All at \textbf{fixed $A$} (process definitions as in
\citealt{langanke2000,langanke2003}; the grouping is verbatim in
\citealt{heger2001} --- EC and $\beta^+$ decay decrease \Ye{}, positron capture
and $\beta^-$ decay increase it). Note the pairing carefully, because it is easy
to get backwards: the processes group by \textbf{which way $Z$ moves}, not by
which lepton is named. EC and $\beta^+$ decay both \emph{lower} $Z$ and \Ye{}
and emit a $\nu$; $\beta^-$ decay and positron \emph{capture} both \emph{raise}
$Z$ and \Ye{} and emit a $\bar\nu$. Grouping ``$\beta^+$ decay / positron
capture'' together --- the intuitive but wrong reading --- puts a \Ye{}-raising
process in the \Ye{}-lowering row, which is exactly the sign error the
project's invariant \#3 exists to catch.

\textbf{What the code carries.} \code{graph/stoich.py::WEAK\_LEDGERS} has three
keys --- \code{electron\_capture}, \code{beta\_neg}, \code{beta\_pos} --- with
EC and $\beta^+$ sharing the ledger signature $(-1, +1, 0)$ under the
annihilated-positron convention (the emitted positron annihilates on a plasma
electron, so the net electron-ledger effect is $-1$). Positron capture has no
separate entry: it is folded into the \emph{evaluated} $\beta^-$-direction weak
rate --- MESA's weaklib sums decay $+$ capture per direction at evaluation time,
although the raw tables do carry a separate positron-capture column before that
sum (and its contribution is small at these conditions ---
\citealt{langanke2000}; \citealt{heger2001} neglected it outright). So the
fourth row above is physics you must know in order to read a weak-rate table
correctly, not a column you will find in \code{weak\_mask}.

Hence a structural statement you should verify against the exported matrix:

\begin{invariant}
A column of $\nu$ is a weak column \textbf{if and only if} it moves \Ye{}, i.e.
$\sum_i Z_i \nu_{ij} \neq 0$ while $\sum_i A_i \nu_{ij} = 0$ --- with
\textbf{both sums taken over the nuclei block only.} On the extended
$\tilde\nu$ (which appends the $e^-$/$\nu$/$\bar\nu$ ledger rows) the charge sum
is zero for \emph{every} column, weak included, because charge is conserved by
the reaction.
\end{invariant}

Verify it, but do not adopt it as a \emph{definition}: \code{graph/stoich.py}
assigns the lepton ledger from the rate's \code{weak\_type} and warns explicitly
against back-deriving it from $Z\cdot\nu$, ``that would make the
charge-to-lepton closure a tautology''. The eligible-mask machinery is built
structurally from \code{weak\_mask} at construction time for the same reason.

\section{Why silicon burning drives \texorpdfstring{\Ye{}}{Ye} down --- now
  quantitatively}
\label{sec:yedown}

Electron capture on free protons requires total electron energy
$E_e \geq (m_n - m_p)c^2 = 1.293$~MeV. In a degenerate gas the available energy
is $E_F$, and setting $E_F = 1.293$~MeV gives a \textbf{critical density
$\rho \approx 2.4\times10^{7}$~\gcc} at $\Ye = 0.5$ (\S\ref{sec:ecthreshold}).

That is \emph{inside the regime box, near its lower edge}. Below it captures are
thermally suppressed; above it they are degeneracy-driven and vigorous. At
$\rho = 10^{8}$ the Fermi energy is 1.97~MeV --- comfortably above threshold ---
and at $10^{9}$ it is 4.1~MeV.

The emitted neutrinos escape (\S\ref{sec:freestreaming}), so nothing accumulates
to push the weak sector back toward equilibrium: \textbf{\Ye{} ratchets downward
in net}, and by \S\ref{sec:explodability} the Chandrasekhar mass follows it. The
individual channels still run both ways --- $\beta^-$ decay raises \Ye{} every
step --- so ``monotonic'' is a statement about the sum over weak columns, not
about any one of them (\S\ref{sec:freestreaming}). That net one-way ratchet over
${\sim}10^{5}$~s is the physical signal this entire project exists to reproduce
accurately.

\section{The lepton ledger}
\label{sec:leptonledger}

Because neutrinos leave, bookkeeping must be explicit. $C$ carries three rows:

\begin{enumerate}
\item \textbf{baryon} --- $A_i$ on nuclei
\item \textbf{charge} --- $Z_i$ on nuclei, \textbf{$-1$ on the electron
      column}, 0 on neutrinos
\item \textbf{lepton number} --- $e^-$ $+1$, $\nu$ $+1$, $\bar\nu$ $-1$
\end{enumerate}

The projector therefore acts on the \textbf{extended} vector
$[\dd Y_{\text{nuclei}}\ldots, \dd Y_{e^-}, \dd Y_\nu, \dd Y_{\bar\nu}]$.

\textbf{Why the extension is not optional.} Projecting a nuclei-only vector
would force $\sum Z_i \dd Y_i = 0$ --- precisely the statement that \Ye{} never
changes. It would silently erase the physical signal. This is the concrete
content of the invariant ``a design that zeroes $\dd\Ye$ to make drift residuals
vanish is wrong'': zeroing the signal is the \emph{easy} way to pass a
conservation test, and the architecture must make it impossible by accident.
\code{test\_projector.py} carries a dedicated
\code{test\_projection\_preserves\_weak\_dYe\_signal} for exactly this.

Note the conceptual subtlety the third row encodes: lepton number is
\emph{conserved by the reaction}, and then the neutrino physically leaves the
zone. The ledger lets you state both facts without contradiction.

% ---------------------------------------------------------------------
\chapter{The network as a bipartite graph}
\label{sec:bipartite}

Two node types:

\begin{itemize}
\item \textbf{Isotope nodes} $I$: one per species (80 or 151).
\item \textbf{Reaction nodes} $R$: one per reaction (607 or 1518 after
      reconciliation).
\end{itemize}

Two edge types:

\begin{itemize}
\item \textbf{$I \to R$} reactant incidence.
\item \textbf{$R \to I$} signed product incidence, carrying the stoichiometric
      coefficient.
\end{itemize}

$\nu$ is exactly the (signed, summed) biadjacency matrix --- the graph and the
linear algebra are two views of one object.

\textbf{Why bipartite rather than species-only.} A reaction with three reactants
is a genuine 3-way interaction --- a \textbf{hyperedge}. Collapsing it to
pairwise $I\to I$ edges loses the information that the three must act
\emph{jointly}: from the projected graph you cannot tell whether $\{a,b,c\}$
came from one ternary reaction or three binary ones, and the rate expressions
differ ($\rho^2 Y_aY_bY_c$ versus sums of $\rho Y\cdot Y$). The bipartite form
represents the hyperedge exactly. The cost is a doubled hop count: an
$I\to I$ step corresponds to two bipartite hops along an alternating path, which
is why the bipartite diameter measures 6 while the derived $I\to I$ view is
measured at diameter 4 (\code{RESULTS.md} --- not exactly half, because the
projection is its own graph object). One message-passing round traverses both
halves --- see \S\ref{sec:messagepassing} for the depth arithmetic and for why
the shipped $K = 5$ is margin rather than a minimum.

\textbf{Forward and reverse are separate columns}, not one signed column. MESA
softwires them separately and pynucastro treats them as distinct rates, so the
canonical key convention makes a directed key per direction, and \code{pair\_col}
recovers the pairing downstream. This is exactly what makes $\kappa$ computable
per pair --- and note that it is a \emph{choice}: a signed-column representation
would make conservation just as exact but would destroy the ability to measure
detailed balance, because $f^+$ and $f^-$ would already have been summed.

% ---------------------------------------------------------------------
\chapter{Reaction nomenclature and REACLIB chapters}
\label{sec:nomenclature}

Basic literacy. The configs index reactions by these categories, and the
project's headline bug is indexed by \emph{chapter}.

\section{Channel notation}
\label{sec:channelnotation}

\begin{equation*}
A(b,c)D \qquad \text{means} \qquad A + b \to c + D
\end{equation*}

with the light participants abbreviated: $n$, $p$, $\alpha$, $\gamma$, $e^-$,
$\nu$ (standard notation --- \citealt{jose2011}). Channels that appear in this
network:

\begin{center}
\small
\begin{tabularx}{\textwidth}{@{}l R{0.864} R{1.136}@{}}
\toprule
\textbf{Channel} & \textbf{Meaning} & \textbf{Role in Si burning} \\
\midrule
$(\alpha,\gamma)$ / $(\gamma,\alpha)$ &
  $\alpha$ capture / $\alpha$ photodisintegration &
  \textbf{the ladder} --- Si $\to$ S $\to$ Ar $\to \ldots \to$ Ni \\
$(p,\gamma)$ / $(\gamma,p)$ & proton capture / release &
  side flow, bottlenecks \\
$(n,\gamma)$ / $(\gamma,n)$ & neutron capture / release &
  fastest pairs; set the stiffness ceiling \\
$(p,\alpha)$, $(\alpha,p)$ & rearrangement & bridges between groups \\
$(n,\alpha)$, $(\alpha,n)$ & rearrangement & bridges \\
$(p,n)$, $(n,p)$ & charge exchange, strong &
  conserves \Ye{} (both $A$ and $Z$ conserved overall) \\
$(e^-,\nu)$ & electron capture & \textbf{weak} --- lowers \Ye{} \\
$\beta^-$, $\beta^+$ & beta decay &
  \textbf{weak} --- raises / lowers \Ye{} \\
\bottomrule
\end{tabularx}
\end{center}

\textbf{A reverse is a distinct reaction}, not a sign flip: $(\alpha,\gamma)$
and $(\gamma,\alpha)$ are two columns of $\nu$, two rate evaluations, two rows
in the config. \S\ref{sec:bipartite} explains why that representational choice
is what makes $\kappa$ measurable.

\section{REACLIB chapters}
\label{sec:reaclibchapters}

REACLIB indexes every rate by a \textbf{chapter}, which is simply the reactant
and product counts \citep[][all eleven forms below verified against the JINA
REACLIB format definition]{cyburt2010}:

\begin{center}
\small
\begin{tabularx}{\textwidth}{@{}c R{0.95} c R{1.05}@{}}
\toprule
\textbf{Chapter} & \textbf{Form} &
\textbf{$\Delta N = n_{\text{out}} - n_{\text{in}}$} & \textbf{Example} \\
\midrule
1  & $e_1 \to e_2$ & 0 & $\beta$ decay, electron capture \\
2  & $e_1 \to e_2 + e_3$ & $+1$ &
     $(\gamma,n)$, $(\gamma,p)$, $(\gamma,\alpha)$ \\
3  & $e_1 \to e_2 + e_3 + e_4$ & $+2$ & $\nuc{C}{12} \to 3\alpha$ \\
4  & $e_1 + e_2 \to e_3$ & $-1$ &
     $(n,\gamma)$, $(p,\gamma)$, $(\alpha,\gamma)$ \\
5  & $e_1 + e_2 \to e_3 + e_4$ & 0 &
     $(p,\alpha)$, $(\alpha,n)$, $(n,p)$ \\
6  & $e_1 + e_2 \to e_3 + e_4 + e_5$ & $+1$ &
     $\nuc{B}{11}(p,2\alpha)\alpha$, $\nuc{He}{3}(\nuc{He}{3},2p)\nuc{He}{4}$
     class \\
7  & $e_1 + e_2 \to e_3 + e_4 + e_5 + e_6$ & $+2$ &
     light-ion $2\to4$ breakup, e.g.\
     $\nuc{Li}{7}(\nuc{He}{3},n p \alpha)\alpha$ \\
8  & $e_1 + e_2 + e_3 \to e_4$ & $-2$ &
     \textbf{$3\alpha \to \nuc{C}{12}$} \\
9  & $e_1 + e_2 + e_3 \to e_4 + e_5$ & $-1$ & \\
10 & $e_1 + e_2 + e_3 + e_4 \to e_5 + e_6$ & $-2$ & \\
11 & $e_1 \to e_2 + e_3 + e_4 + e_5$ & $+3$ & \\
\bottomrule
\end{tabularx}
\end{center}

\textbf{Why the chapter is the quantity that matters.} Detailed balance relates
a forward rate to its reverse through a phase-space factor carrying one power of

\begin{equation*}
\left(\frac{\mu\,kT}{2\pi\hbar^2}\right)^{3/2}\Big/ N_A
\end{equation*}

\textbf{per excess product} --- i.e.\ per unit of $|\Delta N|$
(\S\ref{sec:detailedbalance}, and derived properly in S3; the general
reverse-ratio formula carries the factor to the power $r - p$ ---
\citealt{rauscher2000}; \citealt{smith2023}, appendix). Note the reduced mass
$\mu$ inside: without it the factor is not a number density, and the
bookkeeping does not close dimensionally.

\begin{warn}
\textbf{Do not compress this to ``$|\Delta N| \neq 1$ is the bug''.} That is the
natural guess and it is wrong --- see below. The discriminator is the
\emph{arity} of the tabulated direction, not $|\Delta N|$.
\end{warn}

The gh-575 / Appendix-B bug lives here (\citealp{mesa575}; \citealt[App.~B]{grichener2025}): stock MESA's detailed-balance path handles only the
\textbf{$2 \to 1$} forward (chapter 4, giving the chapter-2 reverse), where
exactly one power of the factor is required, plus the factor-free $2 \to 2$
case. Any tabulated direction with \textbf{three or more participants on a
side} falls outside that path --- the issue reports the phase-space factor is
omitted there, producing reverse rates wrong by up to ${\sim}24$ orders of
magnitude (exactly how many powers are missing per channel is an inference from
the formula, not stated in the issue). So the affected forward chapters are
\textbf{6} ($2\to3$), \textbf{7} ($2\to4$), \textbf{8} ($3\to1$), and
\textbf{9} ($3\to2$), and it is their \emph{inverses} that MESA gets wrong.

\textbf{Read the config with its convention in hand}, because it is not stated
in the file: in \code{configs/appendixb\_excluded\_channels.yaml},
\code{chapter} is the chapter of the \textbf{tabulated forward}, while
\code{mesa\_handle} and \code{dN} describe the \textbf{inverse} --- the
direction actually being flagged. Hence rows like

\begin{figure}[H]
\centering
\begin{tikzpicture}[font=\footnotesize\ttfamily,
  ann/.style={font=\footnotesize\rmfamily\itshape, align=left,
              text=warnrule}]

\node[anchor=west, align=left] (r1) at (0,0)
  {chapter: 8\ \ \ dN: +2\ \ \ \ r\_c12\_to\_he4\_he4\_he4};
\node[anchor=west, align=left] (r2) at (0,-1.35)
  {chapter: 6\ \ \ dN: $-$1\ \ \ \ r\_he4\_he4\_he4\_to\_h1\_b11};

\node[ann, anchor=west] (a1) at (8.35,0.16)
  {forward $3\alpha\to\nuc{C}{12}$ is ch-8;};
\node[ann, anchor=west] (a1b) at (8.35,-0.20)
  {the flagged handle is its $1\to3$ inverse};

\node[ann, anchor=west] (a2) at (8.35,-1.19)
  {forward $h1+b11\to3\alpha$ is ch-6;};
\node[ann, anchor=west] (a2b) at (8.35,-1.55)
  {the flagged handle is its $3\to2$ inverse};

\draw[-{Latex[length=1.6mm]}, warnrule, line width=0.6pt]
  ([xshift=2mm]r1.east) -- ([xshift=-2mm]a1.west);
\draw[-{Latex[length=1.6mm]}, warnrule, line width=0.6pt]
  ([xshift=2mm]r2.east) -- ([xshift=-2mm]a2.west);
\end{tikzpicture}
\caption{Two rows of \code{configs/appendixb\_excluded\_channels.yaml} with
their convention made explicit: \code{chapter} names the tabulated
\emph{forward}, \code{dN} the flagged \emph{inverse}.}
\label{fig:appendixb}
\end{figure}

\begin{exercise}
\textbf{The falsification you should perform yourself.} Count \code{dN} over the
20 rows of that config (11 for \msn{151}, 9 for \msn{80}): \textbf{13 of them
have $|dN| = 1$} --- all eleven chapter-6 rows and both chapter-9 rows. The only
rows with $|dN| \neq 1$ are the two chapter-7 ($dN = -2$) and five chapter-8
($dN = +2$) entries. If ``$|\Delta N| \neq 1$'' were the bug, the other thirteen
rows could not be there. Conversely chapter 8's forward \emph{does} have a
single product, so a rule phrased on product count alone would wrongly exonerate
$3\alpha \to \nuc{C}{12}$ --- whose inverse is in the config with
\code{dlog10} in the tens. The config's own \code{class} values name the real
criterion (\code{multi\_body\_inverse} for the ch-6/7/9 rows,
\code{photo\_of\_multibody} for the five ch-8 rows --- the $3\alpha$ inverse is
the latter). The \code{dlog10} column is how far wrong the rate was.
\end{exercise}

Chapter 5 ($2\to2$) is safe for a different reason again --- $\Delta N = 0$, so
no factor is needed at all. Chapters 10 and 11 do not appear as
tabulated-\emph{forward} chapters: REACLIB's five ch-10 entries are all
v-flagged detailed-balance reverses of the ch-7 forwards (three touch only
in-network species, e.g.\ $n+p+\alpha+\alpha \to \nuc{He}{3}+\nuc{Li}{7}$ ---
i.e.\ exactly the raw-reverse class the pf gate forbids using), and its ch-11
forwards are all $\beta$-delayed multi-neutron emitters near the dripline,
outside these networks. Not because either arity would be safe.

\begin{repopointer}
\textbf{In the code:}
\code{configs/reactions\_mesa\{80,151\}\_mesa.yaml} carries \code{category} per
reaction (see \S\ref{sec:topology});
\code{configs/appendixb\_excluded\_channels.yaml} carries \code{chapter} and
\code{dN}. \code{crosscheck/canonical.py} builds the directed multiset key that
makes ``same reaction'' decidable across MESA and pynucastro.
\end{repopointer}

% ---------------------------------------------------------------------
\chapter{The topology of silicon burning}
\label{sec:topology}

What actually flows where. \S\ref{sec:lastday} said ``photodisintegration
rearrangement''; this is the mechanism in detail.

\section{The \texorpdfstring{$\alpha$}{alpha} ladder}
\label{sec:alphaladder}

Photodisintegration of the least-bound nuclei liberates $\alpha$, $p$, $n$.
Those light particles are captured by the rest, walking composition up an
$\alpha$-chain \citep{bodansky1968,wac1973,hix1996}:

\begin{figure}[H]
\centering
\begin{tikzpicture}[
  font=\small,
  nuc/.style={draw, rounded corners=2pt, line width=0.7pt, inner sep=3.5pt,
              fill=resultrule!7, minimum height=7mm},
  fwd/.style={-{Latex[length=1.7mm]}, line width=0.65pt},
  lbl/.style={font=\scriptsize, inner sep=1.4pt},
  node distance=13mm]

% top row, left to right
\node[nuc] (si)  at (0,0)      {\nuc{Si}{28}};
\node[nuc] (s)   at (2.75,0)   {\nuc{S}{32}};
\node[nuc] (ar)  at (5.5,0)    {\nuc{Ar}{36}};
\node[nuc] (ca)  at (8.25,0)   {\nuc{Ca}{40}};
\node[nuc] (ti)  at (11.0,0)   {\nuc{Ti}{44}};
% bottom row, right to left
\node[nuc] (cr)  at (11.0,-2.3) {\nuc{Cr}{48}};
\node[nuc] (fe)  at (8.25,-2.3) {\nuc{Fe}{52}};
\node[nuc] (ni)  at (5.5,-2.3)  {\nuc{Ni}{56}};

\foreach \a/\b in {si/s, s/ar, ar/ca, ca/ti} {
  \draw[fwd] ([yshift=1.6mm]\a.east) -- ([yshift=1.6mm]\b.west)
    node[lbl, midway, above] {$(\alpha,\gamma)$};
  \draw[fwd] ([yshift=-1.6mm]\b.west) -- ([yshift=-1.6mm]\a.east)
    node[lbl, midway, below] {$(\gamma,\alpha)$};
}

% Ti44 -> Cr48 turn
\draw[fwd] (ti.south) -- (cr.north);

\foreach \a/\b in {cr/fe, fe/ni} {
  \draw[fwd] (\a.west) -- (\b.east)
    node[lbl, midway, above] {$(\alpha,\gamma)$};
}
\end{tikzpicture}
\caption{The $\alpha$ ladder. Each rung is a capture/photodisintegration pair;
the chain turns at \nuc{Ti}{44} and runs back along the lower row to the
\nuc{Ni}{56} endpoint.}
\label{fig:alphaladder}
\end{figure}

\textbf{The config category names are literally this ladder.} Count them in
\code{configs/reactions\_mesa80\_mesa.yaml}:

\begin{center}
\small
\begin{tabularx}{\textwidth}{@{}R{0.955} l R{1.045}@{}}
\toprule
\textbf{Category} & \textbf{Count (\msn{80})} & \textbf{What it is} \\
\midrule
\code{photo} & \textbf{151} &
  photodisintegration --- the largest single category \\
\code{si\_alpha}, \code{s\_alpha}, \code{ar\_alpha}, \code{ca\_alpha},
\code{ti\_alpha}, \code{cr\_alpha} & 27, 22, 19, 17, 14, 14 &
  the rungs of the ladder \\
\code{fe\_co\_ni} & 17 & the top of the ladder \\
\code{c\_alpha}, \code{n\_alpha}, \code{o\_alpha}, \code{ne\_alpha},
\code{na\_alpha}, \code{mg\_alpha} & 7, 10, 23, 32, 25, 27 &
  below Si --- inherited from earlier stages \\
\code{tri\_alpha} & 1 &
  $3\alpha \to \nuc{C}{12}$ (its $1\to3$ inverse is the gh-575 casualty,
  \S\ref{sec:reaclibchapters}) \\
\code{c12\_c12}, \code{o16\_o16} & 3, 3 &
  heavy-ion fusion, earlier stages \\
\code{pp}, \code{cno} & 15, 33 &
  hydrogen burning, essentially inert here \\
\code{other} & 147 & $(p,n)$, $(n,p)$, weak, everything else \\
\midrule
\textbf{total} & \textbf{607} & $\leftarrow$ the table sums; check it \\
\bottomrule
\end{tabularx}
\end{center}

That \code{photo} is the largest category is the whole story of
\S\ref{sec:lastday} in one number.

\section{Two clusters and the bridges between them}
\label{sec:clusters}

Both ends of the ladder equilibrate internally faster than material crosses the
middle. This produces the structure the project's QSE machinery targets --- two
QSE groups divided near $A \approx 45$ \citep{wac1973,hix1996}:

\begin{figure}[H]
\centering
\begin{tikzpicture}[font=\small,
  grp/.style={draw, line width=0.8pt, rounded corners=3pt, align=center,
              minimum width=46mm, minimum height=20mm, fill=resultrule!6},
  lbl/.style={font=\scriptsize, align=center}]

\node[grp] (si) at (0,0)
  {\textbf{SILICON GROUP}\\[2pt] $24 \le A < 45$\\[2pt]
   \emph{internally fast}};
\node[grp] (fe) at (8.4,0)
  {\textbf{IRON GROUP}\\[2pt] $A \gtrsim 45$\\[2pt]
   \emph{internally fast}};

\draw[-{Latex[length=2.4mm]}, line width=1.1pt]
  (si.east) -- (fe.west)
  node[lbl, midway, above=1pt] {bridges}
  node[lbl, midway, below=1pt] {bottleneck\\reactions};

\draw[-{Latex[length=2.2mm]}, line width=0.7pt]
  (fe.south) -- ++(0,-1.1) coordinate (bend)
  -| (si.south);
\node[lbl, fill=white, inner sep=2pt] at (4.2,-1.85)
  {free $n$, $p$, $\alpha$};
\end{tikzpicture}
\caption{The two QSE clusters, each internally equilibrated, coupled by a
handful of rate-limiting bridge reactions and by the shared reservoir of free
light particles.}
\label{fig:clusters}
\end{figure}

\begin{itemize}
\item \textbf{Within a group}, forward and reverse rates are fast and nearly
      balanced --- so $\kappa$ is small there, and the abundances follow an
      equilibrium (Saha-like) distribution up to one cluster-wide offset
      (\S\ref{sec:chempotential}, eq.~\eqref{eq:qse-offset}).
\item \textbf{Between groups}, flow is rate-limited by a handful of
      \textbf{bridge} reactions. These carry the net baryon flux and therefore
      control how fast Si becomes Fe.
\end{itemize}

This is what \code{configs/qse\_groups.yaml} encodes (\code{a\_min: 24},
\code{a\_max: 45}, with the \code{a24\_46} variant moving the boundary so that
\nuc{Sc}{45} falls inside), and what the bridge analysis of S9 measures. It also
explains why a \emph{bottleneck} reaction like
$\nuc{Sc}{45}(p,\gamma)\nuc{Ti}{46}$ is worth naming individually: a small
number of reactions carry a large share of the inter-group flow --- in the
classic analysis that single reaction dominated the linkage, with perhaps a
quarter of the flow through lesser channels \citep[][as quoted by
\citealt{hix1996}]{wac1973}.

\section{A note on energy accounting}
\label{sec:energyaccounting}

Two independent routes to the same specific energy release, both used by the
project as a mutual check:

\begin{align*}
\text{composition route:}\quad \enuc &= -N_A\sum_i m_i c^2\,\Delta Y_i
  && \text{(mass excesses, \S\ref{sec:massexcess})}\\
\text{flux route:}\quad \enuc &= \sum_j Q_j\,\Phi_j
  && \text{($Q$-values} \times \text{fluxes)}
\end{align*}

They must agree, because $Q$-values are themselves differences of mass excesses.
Agreement is therefore \emph{near-algebraic} when $Q$ is taken constant --- it
bounds $Q$-table rounding, not route physics. A genuinely independent check
needs the energy from one code's rates against another's composition change,
which is why the project reports two separate readings and treats only the
second as informative. Keep that distinction; it is easy to over-claim here.

% ---------------------------------------------------------------------
\chapter{The timescale hierarchy}
\label{sec:timescales}

Every structural feature of this problem is a statement about a ratio of
timescales. Assembling them in one place is the single most useful organizing
device in Tier 0.

\section{The species destruction timescale}
\label{sec:tauspecies}

For species $i$, define

\begin{equation}
\tau_i = \frac{Y_i}{|\dot{Y}_i|}
\label{eq:tau-species}
\end{equation}

--- the e-folding time of that species at the current state. This is the
quantity that decides whether a label step is ``linear'' ($\tau \gg \Delta t$)
or a full relaxation ($\tau \lesssim \Delta t$), and it is what the handshake
analysis bins on (\S\ref{sec:dtgrid}).

\section{The hierarchy, at box centre}
\label{sec:hierarchy}

Order-of-magnitude estimates at $\Tn \approx 4$, $\rho \approx 10^{8}$ ---
derive each yourself rather than trusting the table:

\begin{center}
\small
\begin{tabularx}{\textwidth}{@{}R{0.930} l R{1.070}@{}}
\toprule
\textbf{Process} & \textbf{Timescale} & \textbf{Set by} \\
\midrule
$(\gamma,n)$ / $(n,\gamma)$ on loosely bound nuclei &
  $10^{-10}$--$10^{-6}$~s & $e^{-Q/kT}$ with $Q \sim$ few MeV \\
charged-particle captures & $10^{-6}$--$10^{0}$~s &
  Coulomb barrier $+$ screening \\
intra-group equilibration & $\lesssim 10^{-3}$~s &
  the above, within a cluster \\
\textbf{inter-group (bridge) flow} & $10^{-1}$--$10^{2}$~s &
  bottleneck reactions \\
\textbf{weak reactions (EC, $\beta$)} & $10^{2}$--$10^{5}$~s &
  $ft$ values, $E_F$ \\
label $\dd t$ grid & $10^{-6}$--$10^{2}$~s &
  $\leftarrow$ the emulator's window \\
convective turnover & $10^{1}$--$10^{3}$~s &
  $v_{\mathrm{conv}} \sim 10^{6}$--$10^{7}$~cm/s over
  $L \sim 10^{8}$--$10^{9}$~cm \\
core Si burning duration &
  ${\sim}10^{5}$--$10^{6}$~s ($\approx$1 day to ${\sim}$2 weeks,
  mass-dependent; WHW02) & neutrino cooling \\
free-fall collapse & ${\sim}0.07$~s &
  $\sqrt{3\pi/32G\rho}$ at $\rho = 10^{9}$ \\
\bottomrule
\end{tabularx}
\end{center}

\textbf{Read the consequences straight off the table:}

\begin{itemize}
\item \textbf{Stiffness ratio} (\S\ref{sec:stiffness}) is the ratio of the
      extremes: ${\sim}10^{10}/10^{-5} \approx 10^{15}$.
\item \textbf{QSE exists} because intra-group $\ll$ inter-group. The separation
      \emph{is} the cluster structure.
\item \textbf{\Ye{} is slow} --- weak timescales are the longest in the burning
      problem, which is why \Ye{} evolution is the thing you must integrate
      accurately for a long time, and why it has no restoring force on the fast
      manifold (\S\ref{sec:rollouterror}).
\item \textbf{The $\dd t$ grid brackets the interesting range} --- its
      $10^{-6}$~s floor sits at (not below) the fastest charged-particle capture
      timescales, and its top end reaches past the bridge timescale.
\item \textbf{Operator splitting is safe at the short end, questionable at the
      long end} (\S\ref{sec:splitting}): $\Delta t = 10^{2}$~s is comparable to
      the bridge and convective timescales.
\item \textbf{One-zone is an approximation} because convective turnover
      ($10^{1}$--$10^{3}$~s) is \emph{not} long compared to the bridge timescale
      (\S\ref{sec:onezone}).
\end{itemize}

\begin{warn}
\textbf{Do the convective row's arithmetic yourself, and notice it is a range
you cannot pin.} $\tau_{\mathrm{conv}} = L/v_{\mathrm{conv}}$, and both factors
are uncertain by a decade: $L \sim 10^{8}$--$10^{9}$~cm for the convective
Si-burning core, $v_{\mathrm{conv}} \sim 10^{6}$--$10^{7}$~cm/s. The corners
give 10~s and $10^{3}$~s. Quoting a tight ``$10^{2}$--$10^{3}$~s'' requires
committing to $L \sim 10^{9}$~cm, which is defensible but is a choice, not a
derivation. This matters because two conclusions lean on the comparison:
whether $\Delta t = 10^{2}$~s violates the splitting assumption
(\S\ref{sec:splitting}) and whether one-zone is safe (\S\ref{sec:onezone}). At
the fast corner mixing is \emph{faster} than the bridge flow and one-zone is
worse than the table suggests; at the slow corner they are comparable. Carry
the range, not a point value.
\end{warn}

\section{The separation you are tempted to assume does not exist}
\label{sec:nogap}

The natural next move is: ``the fast sector is 6--8 orders faster than the slow
sector, so equilibrate it algebraically and integrate only the slow part.''

That figure was carried in this project as an \emph{assumed} number and has been
\textbf{retired by measurement}. On the actual label manifold, per-stratum
medians of

\begin{equation*}
\log_{10}\left[\frac{\text{fastest } \kappa\text{-balanced gross rate}}
                    {\text{slowest inter-group bottleneck net rate}}\right]
\end{equation*}

span roughly \textbf{$-12$ to $+1.4$~dex} ($p_{10} \approx -25$,
$p_{90} \approx +3.5$). There is no cleanly separated fast equilibrated sector:
$\kappa$-balanced pairs are $\le 0.4\%$ of carrying pairs and are mostly
low-flux.

\begin{aside}
\textbf{So the hierarchy above is a real and useful organizing device, but it is
a hierarchy of \emph{typical} values with enormous overlap --- not a spectral
gap.} Any design that assumes a clean gap (a frozen equilibrium mask, a
two-timescale reduction) has to justify itself against that measurement. This is
the quantitative reason the Guidry mask came up empty and the flux head ended up
full-width.
\end{aside}

% ---------------------------------------------------------------------
\chapter{Code walkthrough}
\label{sec:codewalkthrough-s0}

\textbf{\code{configs/isotopes\_mesa80.yaml} / \code{isotopes\_mesa151.yaml}}
--- authoritative species tables. Read the provenance headers: lists come from
the \emph{data column headers}, cross-checked against the published paper's
appendix table, which \textbf{omits \nuc{Ca}{41}}. Resolution recorded: ``paper
typo; data is authoritative.'' The working principle it establishes --- when
paper and data disagree, the data wins, and the disagreement is recorded rather
than smoothed over --- recurs throughout the project.

\textbf{\code{graph/isotopes.py}} --- \code{IsotopeTable},
\code{load\_isotope\_table}, \code{canonical\_network}. The job is naming
discipline: YAML names and order define $\nu$'s row order, the npz
\code{species} array, and graph node names, matching the training-CSV column
suffixes \code{initial\_<name>}. pynucastro's capitalization (\code{"Neut"},
\code{"H1"}) is \textbf{never stored} --- mapping is fail-loud at the boundary.

\begin{repopointer}
\textbf{Oracle:} \code{tests/test\_schema.py::test\_network\_species\_counts}\\
\textbf{Notebook:} \code{01-data-inventory} Fig~2 (($N$,$Z$) plane), Fig~3
(\Ye{}-controller membership)
\end{repopointer}

% ---------------------------------------------------------------------
\chapter*{Self-check --- the network as a formal object}
\addcontentsline{toc}{chapter}{Self-check --- the network as a formal object}
\label{sec:selfcheck-network}

\begin{selfcheck}
\item Derive $Y = X/A$ and $\sum A_i Y_i = 1$ from
      $n_i = \rho N_A X_i/A_i$ without looking.
\item Show that the \emph{only} columns of $\nu$ with
      $\sum Z_i \nu_{ij} \neq 0$ are weak columns, and confirm numerically
      against the exported \code{weak\_mask}.
\item Given $M_{\mathrm{ch}} = 5.83\,Y_e^2\,\Msun$, compute the fractional
      change for $\Ye: 0.50 \to 0.47 \to 0.45$. Then compute what per-step \Ye{}
      error, accumulated systematically over $1.6\times10^{3}$ steps, produces a
      1\% $M_{\mathrm{ch}}$ error. Compare to the $3\times10^{-6}$ gate.
\item Explain why \code{neut} must be a network species rather than a bath, in
      terms of equation \eqref{eq:baryon-conservation}.
\item Explain the \S\ref{sec:molarabundance} precision note: why use integer $A$
      in the constraint when $M \neq A\cdot u$? Compute $(M - A u)/(A u)$
      yourself for \code{neut}, \code{he4}, and \code{fe56} from the mass
      excesses in \S\ref{sec:massexcess} --- get the signs right --- and say
      which of those numbers sets the $10^{-3}$ that $A\cdot\nu = 0$ would
      otherwise be limited to. What breaks if you use real masses in the
      constraint row instead?
\item Explain what is lost by projecting the bipartite graph onto species-only
      edges, using a concrete ternary reaction from the network.
\item Write $\nuc{Ca}{40}(\alpha,\gamma)\nuc{Ti}{44}$ and its reverse in chapter
      notation. What is $\Delta N$ for each? Why is this pair safe --- and state
      the reason in terms of arity, not $|\Delta N|$.
\item Take three entries from
      \code{configs/appendixb\_excluded\_channels.yaml} and verify both fields
      against the table in \S\ref{sec:reaclibchapters} --- remembering that
      \code{chapter} describes the tabulated \emph{forward} while \code{dN}
      describes the flagged \emph{inverse}. Then tabulate $|dN|$ over all 20
      rows and confirm that most of them are 1, i.e.\ that
      $|\Delta N| \neq 1$ is \emph{not} the criterion.
\item Why is \code{photo} the largest reaction category in the network? Connect
      to \S\ref{sec:lastday}.
\item Estimate $\tau$ for a $(\gamma,n)$ with $Q = 8$~MeV at $\Tn = 4$ and at
      $\Tn = 5$. How many decades does one GK buy you?
\item Given the hierarchy table, predict which pairs of processes could
      plausibly \emph{swap order} somewhere in the box. Check against the
      $-12\ldots+1.4$~dex measurement.
\item State in one sentence why ``6--8 orders of separation'' being false does
      not invalidate the QSE \emph{concept}, only its algebraic exploitation.
\end{selfcheck}

\begin{aside}
\textbf{Two questions from the source document's \S II.8 are not here.} Item~4
(the \msn{80}/\msn{151} diff and its \Ye{} bias) and item~7 (the $\asinh$
constraint counter-example) both test material that now sits later, in
Part~\ref{part:computational}; they appear in that part's self-check instead.
See Appendix~\ref{app:concordance}.
\end{aside}
